\section{Joint analyticity at the grazing boundary}
\label{sec:joint}

The boundary is now analytic, but the given foliation is still only
continuous.  We next construct a parametrization that is holomorphic in
both its circular variable and its translation parameter.  The principal
issue is division by difference operators at rational parameters.  The
periodic graphs from Proposition~\ref{prop:graph-collar} provide two
compatibility conditions: one for the residual and one for its first
parameter derivative.  Both are needed.

We follow the Newton construction of \cite[Section~3]{NB}, giving the
arguments that identify its geometric hypotheses in the present setting.
The only geometric input in this Section is the invariant graph collar
and the stationary-chain closure of Lemma~\ref{lem:stationary-closure}.
The Fourier division estimate will be stated separately in its precise
form.

\begin{theorem}[Joint analytic parametrization]
\label{thm:joint}
There are positive numbers \(s,d\), with \(d<s\), and a function
\(V(y,t)\), holomorphic on \(\{|\Im y|<2s,\ |t|<d\}\), such that
\(V-y\) is \(2\pi\)-periodic in \(y\), \(V\) is even in \(t\), and
\[
 V(\bar y,\bar t)=\overline{V(y,t)},\qquad V(y,0)=y.
\]
Moreover \(V_y\) is uniformly close to one, and
\begin{equation}
 D_1\bigl(V(y,t),V(y+t,t)\bigr)
 +D_2\bigl(V(y-t,t),V(y,t)\bigr)=0
 \label{joint:stationary}
\end{equation}
for \(|\Im y|<s\), \(|t|<d\).  In particular \(V_y>0\) on the
real domain.
\end{theorem}

\subsection{The normalized equation}

The analytic action obtained in Section~\ref{sec:analytic} is
antisymmetric and diagonally \(2\pi\)-periodic.  On a fixed complex
neighborhood of the real diagonal it satisfies
\begin{equation}
 D(x,x+b)=-\frac{b^3}{24}+O(b^5),\qquad
 D_1(x,x+b)=\frac{b^2}{8}+b^4B(x,b),
 \label{joint:action}
\end{equation}
where \(B\) is holomorphic and periodic in \(x\).  All domains below
are contained in this fixed action neighborhood.

Write \(T_zf(y,t)=f(y+z,t)\), use a prime for a \(y\)-derivative,
and set
\[
 \partial f=\frac{T_{t/2}f-T_{-t/2}f}{t}.
\]
This symbol denotes a difference operator, not a parameter derivative.
For a candidate \(V\), define
\begin{align}
 \mathcal R(V)&=D_1(V,T_tV)+D_2(T_{-t}V,V),\notag\\
 E(V)&=t^{-3}V'\mathcal R(V),
 \label{joint:residual}\\
 a(V)&=t^{-1}D_{12}(T_{-t/2}V,T_{t/2}V)
                  (T_{-t/2}V')(T_{t/2}V').
 \label{joint:a}
\end{align}
For a periodic function, \(\langle f\rangle\) denotes its mean in \(y\).
An assigned first parameter jet prescribes a parameter derivative at one
fixed value of \(t\), without assuming a nearby family of exact solutions.

\begin{lemma}[Regularity and factorization]
\label{joint:factorization}
The expressions \(E\) and \(a\) extend holomorphically across \(t=0\),
with
\begin{equation}
 E(V)(y,0)=\frac{(V')^2V''}{4},\qquad
 a(V)(y,0)=\frac{(V')^3}{4}.
 \label{joint:at-zero}
\end{equation}
They are even in \(t\) whenever \(V\) is even, and
\(\langle E(V)\rangle=0\).  The linearization in a variation \(V'w\)
is
\begin{equation}
 E_V(V'w)=\partial(a\partial w)+(Ew)'.
 \label{joint:linearization}
\end{equation}
At an exact solution at a nonzero real parameter \(t_0\), an assigned
parameter jet \(\dot V=V'z\) gives the total residual jet
\begin{equation}
 \dot E=\partial\left(a\partial z+\frac{a}{t_0}\right).
 \label{joint:parameter-jet}
\end{equation}
Here the dot differentiates all parameter dependence, including the
explicit shifts.
\end{lemma}

\begin{proof}
Antisymmetry gives
\(\mathcal R=D_1(V,T_tV)-D_1(V,T_{-t}V)\).  Put
\[
 A_\pm=\frac{T_{\pm t}V-V}{\pm t}
       =\int_0^1 V'(y\pm ut,t)\,\dd u.
\]
By \eqref{joint:action},
\begin{equation}
 \frac{E}{V'}=
 \frac{A_++A_-}{8}\frac{A_+-A_-}{t}
 +t\bigl[A_+^4B(V,tA_+)-A_-^4B(V,-tA_-)\bigr].
 \label{joint:regular-formula}
\end{equation}
The remaining divided difference has the nonsingular expression
\[
 \frac{A_+-A_-}{t}
 =\int_0^1\int_{-u}^{u}V''(y+vt,t)\,\dd v\,\dd u.
\]
These formulas prove regularity and the first identity in
\eqref{joint:at-zero}.  Since
\(D_{12}(x,x+b)=b/4+O(b^3)\), the second follows as well.  They also
show that the first two variations of \(E\), on domains with room for
the shifts, have bounds in terms of a fixed number of derivatives
through order two in \(y\).  No negative power of \(t\) is needed in
these bounds.  Antisymmetry gives the stated parity.

For real \(t\), the mean action
\[
 \frac1{2\pi}\int_0^{2\pi}
 D\bigl(V(y+\eta,t),V(y+\eta+t,t)\bigr)\,\dd y
\]
is independent of \(\eta\).  Differentiating at zero and shifting one
integral gives \(\langle V'\mathcal R\rangle=0\).  The mean-zero
identity for complex \(t\) follows by holomorphic continuation.

For \eqref{joint:linearization}, separate each off-diagonal variation
into its value at \(y\) and its difference from that value.  With
\[
 c_+=D_{12}(V,T_tV)V'T_tV',\qquad
 c_-=D_{12}(T_{-t}V,V)T_{-t}V'V',
\]
the difference terms are
\[
 t^{-3}\bigl[c_+(T_tw-w)+c_-(T_{-t}w-w)\bigr]
 =\partial(a\partial w),
\]
since \(T_{\pm t/2}a=c_\pm/t\).  The common-value terms, including
the variation of \(V'\), are \((Ew)'\).

At an exact solution, differentiating the factor \(t^{-3}\) contributes
zero.  Differentiation of the explicit shifts contributes
\((c_+-c_-)/t_0^3=\partial a/t_0\), whereas the assigned jet contributes
\(\partial(a\partial z)\).  This proves
\eqref{joint:parameter-jet}.
\end{proof}

We will use candidates with \(V'\) close to one and \(|t|\) small,
so \eqref{joint:a} makes \(a\) close to \(1/4\).  For a Newton
correction of the form \(\Delta V=V'v\), the main
equation is \(\partial(a\partial v)\approx-E\): the remaining
linearization term \((Ev)'\) is quadratic in the residual and the
correction once \(v\) is controlled by the residual.  We will solve
this equation by two approximate primitives,
\[
 \partial H\approx-E,\qquad
 b=a^{-1}(H+\lambda(t)),\quad \langle b\rangle=0,
 \qquad \partial v\approx b.
\]
The free constant \(\lambda(t)\) makes the second right-hand side
have zero mean; each approximation will have error quadratic in the
residual norm.  At a rational parameter, the Fourier multiplier of
\(\partial\) vanishes in some nonzero modes.  Periodic orbits give
compatible residual values for the first division.  At those modes,
the value of the resulting primitive is determined by a parameter
derivative, so compatible first parameter jets are needed to justify
the second division.  Both compatibilities will hold up to quadratic
errors, using jets assigned at individual rational parameters without
differentiating the given continuous foliation.  We first establish
these rational compatibilities and then carry out both divisions with
uniform bounds.

\subsection{Domains and the Newton estimate}

For \(\nu\in\{0,\tfrac12,1\}\), set
\begin{equation}
 \mathcal D^\nu(s,d)=
 \{(y,t):|t|<d,\ |\Im y|<s-\nu|\Im t|\}.
 \label{joint:domains}
\end{equation}
These three domains track the complex shifts in the two divisions.
Starting with \(V\) on \(\mathcal D^0\), the shifts in the residual
require \(\mathcal D^1\).  The first Fourier division will recover
one half of this loss, giving \(H\) on \(\mathcal D^{1/2}\); the
second will give \(v\) on \(\mathcal D^0\).  The radii will decrease
slightly at each operation to control derivatives.

Choose \(d<\min(1,s/2)\) small.  Besides staying in the fixed action domain,
we require that every positive real step in the normalized coordinate
of size at most \(3d\) gives a line in the graph collar, and that all
positive rational translations less than \(d\) lie in the range of
Proposition~\ref{prop:graph-collar}.  Both requirements hold after
shrinking \(d\): the line momentum tends uniformly to the grazing
support as the step tends to zero.
We also choose \(3d\sup_{x\in\R}\phi'(x)<\kappa\), where
\(\kappa\) is the angular-step bound in
Lemma~\ref{lem:stationary-closure}.  Thus every such normalized step
has angular increment less than \(\kappa\).

Candidates are holomorphic on \(\mathcal D^0(s,d)\), real-symmetric,
even in \(t\), periodic modulo \(y\), and satisfy a fixed small bound
on \(V-y\) and its first three \(y\)-derivatives.  In particular
\(V'\) is close to one.  Equation~\eqref{joint:a} and
\eqref{joint:action} then make \(a\) uniformly close to \(1/4\) on
\(\mathcal D^{1/2}(s,d)\).

For \(0<\delta<d/100\), abbreviate
\[
 \mathcal D_j^\nu=\mathcal D^\nu(s-j\delta,d-j\delta),
 \qquad Y_j=\{y:|\Im y|<s-j\delta\}.
\]
Norms without a derivative subscript are supremum norms.  Fixed Cauchy
losses between these domains cost fixed powers of \(\delta^{-1}\).
Indeed polydisks of radius \(\delta/4\) about points of
\(\mathcal D_{j+1}^\nu\) lie in \(\mathcal D_j^\nu\).  A shift by
\(t/2\) increases the required value of \(\nu\) by \(1/2\).

\begin{proposition}[Uniform Newton step]
\label{joint:newton}
There are fixed constants \(C_0,B_0\), depending only on the fixed
analytic action bounds and candidate smallness bounds, for which the
following holds.  Suppose
\[
 \|E(V)\|_{\mathcal D_0^1}\le\eps<1,
 \qquad e^{-N\delta}\le\eps,
 \qquad Q=1+N+\delta^{-1},
\]
where \(N\ge1\) is an integer.  If \(C_0Q^{B_0}\eps\) is sufficiently
small, there is a periodic, real-symmetric, even correction \(\Delta V\)
such that
\begin{align}
 \max_{0\le j\le3}\|\partial_y^j\Delta V\|_{\mathcal D_{20}^0}
   &\le C_0Q^{B_0}\eps,
   \label{joint:step-size}\\
 \|E(V+\Delta V)\|_{\mathcal D_{20}^1}
   &\le C_0Q^{B_0}\eps^2.
   \label{joint:step-error}
\end{align}
In particular \(B_0\) is independent of the radius \(d\), the retained
rational parameters, and the order of a formal initializer.
\end{proposition}

We prove this version of \cite[Proposition~3.3]{NB} from the graph
collar.  In its proof, \(O_{\mathrm p}(\eps^k)\) means a bound
\(CQ^B\eps^k\), where \(C,B\) are fixed.  These constants may increase
a fixed finite number of times.  At the end, a single choice of
\(C_0,B_0\) includes every bound and every polynomial smallness
requirement.  None of the following operations increases the number of
derivatives with \(N\) or with the initializer order.

\subsection{Rational values and first parameter jets}

\begin{lemma}[Rational compatibility]
\label{joint:compatibility}
Let \(t_0=2\pi l/q\ne0\) be reduced, with \(q\le N\) and
\(|t_0|<d-2\delta\).  Under the hypotheses and polynomial smallness
conditions of Proposition~\ref{joint:newton}, there are holomorphic
periodic jets \(w,\dot w\) on \(Y_{10}\), real-symmetric in \(y\),
satisfying
\begin{align}
 w,\dot w&=O_{\mathrm p}(\eps),\notag\\
 E+\partial(a\partial w)&=O_{\mathrm p}(\eps^2),
 \label{joint:value-compatibility}\\
 \bigl(E+\partial(a\partial w)\bigr)^{\displaystyle\cdot}
   &=O_{\mathrm p}(\eps^2).
 \label{joint:jet-compatibility}
\end{align}
The dot in the last line uses the assigned jet \(\dot w\) and
differentiates every occurrence of \(t\).
\end{lemma}

\begin{proof}
We begin with \(t_0>0\).  Write \(U_j(y)=V(y+jt_0,t_0)\) and seek
interior-stationary nodes
\[
 X_j=U_j+U'_j\xi_j,\qquad 0\le j\le q,
 \qquad \xi_0=\xi_q=0.
\]
The endpoint increment is \(2\pi l\).  Multiply the interior
stationarity equation at node \(j\) by the fixed factor
\(t_0^{-3}U'_j\).  At \(\xi=0\) the residual is
\(E(y+jt_0,t_0)\).  By \eqref{joint:linearization}, its derivative in
\(\xi\) is the Dirichlet flux matrix
\begin{equation}
 (A_0\xi)_j=t_0^{-2}
 \bigl[a_{j+1/2}(\xi_{j+1}-\xi_j)
       -a_{j-1/2}(\xi_j-\xi_{j-1})\bigr],
 \quad a_{j+1/2}=a(y+(j+1/2)t_0,t_0),
 \label{joint:dirichlet}
\end{equation}
plus the diagonal potential
\[
 W_j=T_{jt_0}\left(E'-\frac{V''}{V'}E\right).
\]
The prefactor in the finite-dimensional equations is frozen; subtracting
its variation explains the second term in this potential.

The inverse of \(A_0\) has a uniform complex bound.  For
\(A_0\xi=f\), let
\(J_{j+1/2}=a_{j+1/2}(\xi_{j+1}-\xi_j)\).  Then
\[
 J_{j+1/2}=J_{1/2}+t_0^2\sum_{i=1}^j f_i.
\]
The condition \(\sum_{j=0}^{q-1}(\xi_{j+1}-\xi_j)=0\) determines
\(J_{1/2}\).  Since each \(a_{j+1/2}^{-1}\) is close to four, the
modulus of their sum is at least a fixed positive multiple of \(q\).
Summing the resulting increments gives
\begin{equation}
 \|A_0^{-1}\|_{\infty\to\infty}\le Cq^2t_0^2\le CQ^2.
 \label{joint:dirichlet-inverse}
\end{equation}
On \(Y_1\), the potential is bounded by \(CQ\eps\).  Thus a Neumann
series gives the same bound for the full inverse if \(CQ^3\eps\) is
small.

Reducedness supplies the useful lower bound
\begin{equation}
 |t_0|^{-1}\le q/(2\pi)\le N/(2\pi).
 \label{joint:rational-bound}
\end{equation}
The nonlinear equations involve only three neighboring nodes, bounded
fixed derivatives of \(D\), and the factor \(t_0^{-3}\).  Their Hessian
in the supremum norm therefore has size at most \(CQ^3\), independently
of the number of nodes.  Applying the linearized inverse to the
nonlinear remainder defines a contraction on a ball of radius
\(CQ^2\eps\) when \(CQ^7\eps\) is small.  Hence
\(\xi=O_{\mathrm p}(\eps)\) on \(Y_1\), holomorphically in \(y\).
Uniqueness preserves periodicity and real symmetry.  Fixed derivatives
are controlled on smaller strips by Cauchy's estimate.

For real \(y\), the uncorrected increments are bounded below by a fixed
multiple of \(t_0\).  The correction, using
\eqref{joint:rational-bound} and a stronger fixed polynomial smallness
condition, preserves positivity and keeps every step below \(3d\).
After applying \(\phi\), these are actual stationary billiard chains:
subtracting the endpoint primitive in the definition of \(D\) telescopes,
and the coordinate change has nonzero derivative.  Their initial
states lie in the graph collar.  Lemma~\ref{lem:stationary-closure}
therefore gives momentum return as well as angle return.  Each chain
is periodic with rotation \(l/q\), and
Proposition~\ref{prop:graph-collar} identifies the same unique invariant
graph for every starting \(y\).

Invert \(V(\,\cdot\,,t_0)\) locally around the uncorrected nodes to
write
\[
 X_j(y)=V(\widetilde y_j(y),t_0),\qquad
 \widetilde y_j=y+jt_0+O_{\mathrm p}(\eps)
\]
on \(Y_2\).  This inverse is obtained by contraction for the small
displacement from \(y+jt_0\); it requires no assertion about a global
complex inverse.  On the real axis, uniqueness of the rational graph
implies
\(\widetilde y_j=\widetilde y_1^{\,j}\) and
\(\widetilde y_q=y+2\pi l\).  Consequently
\[
 G(y)=\frac1q\sum_{j=0}^{q-1}\bigl(\widetilde y_j(y)-jt_0\bigr)
 \quad\hbox{satisfies}\quad
 G(\widetilde y_1(y))=G(y)+t_0.
\]
Its displacement and, on \(Y_3\), its derivative displacement are
\(O_{\mathrm p}(\eps)\).  Contraction gives an inverse on \(Y_4\).
Thus
\begin{equation}
 V_*(y)=V(G^{-1}(y),t_0)
 \label{joint:rational-solution}
\end{equation}
is an exact solution at \(t_0\), with
\(V_*-V(\,\cdot\,,t_0)=O_{\mathrm p}(\eps)\).  The stationarity
identity follows first on the real axis from the graph dynamics and
then throughout the strip by holomorphy.

We must also give \(V_*\) an appropriate first parameter jet.  Initially
assign \(\dot V_*=V_t(\,\cdot\,,t_0)\).  Its residual jet is
\(O_{\mathrm p}(\eps)\) on \(Y_6\).  Indeed Cauchy's estimate gives
\(E_t=O_{\mathrm p}(\eps)\) there: the distance of \(t_0\) from the
parameter boundary exceeds \(2\delta\).  A polydisk of radius
\(\delta/4\) in each variable about \((y,t_0)\), with \(y\in Y_6\),
lies in \(\mathcal D_0^1\), so this estimate costs only \(C\delta^{-1}\).
Replacing \(V\) by \(V_*\), with the same assigned
jet, changes that derivative by \(O_{\mathrm p}(\eps)\).
To make its polynomial character explicit, if \(J\) is any assigned
jet, direct differentiation gives
\begin{align}
 \dot E={}&-3t^{-4}V'\mathcal R+t^{-3}J'\mathcal R\notag\\
 &+t^{-3}V'\bigl[
 D_{11}(V,T_tV)J+D_{12}(V,T_tV)(T_tJ+T_tV')\notag\\
 &\hspace{29mm}
 +D_{21}(T_{-t}V,V)(T_{-t}J-T_{-t}V')
 +D_{22}(T_{-t}V,V)J\bigr].
 \label{joint:raw-jet}
\end{align}
Only fixed derivatives occur, \(V_t\) and its needed \(y\)-derivatives
have Cauchy bounds, and \eqref{joint:rational-bound} controls every
denominator.

Let \(\Pi=q^{-1}\sum_{j=0}^{q-1}T_{jt_0}\).  This is a projection
on periodic functions, and \(\Pi\partial=\partial\Pi=0\).
Equation~\eqref{joint:parameter-jet} at the exact solution shows that
\(\Pi\dot E_*=0\).  Put \(a_*=a(V_*)\); after the fixed derivative
losses, it is still close to \(1/4\).  On \(\ker\Pi\), the operator
\(\partial(a_*\partial)\) has an inverse with polynomial norm.  For
an explicit verification, put \(T=T_{t_0}\) and
\[
 R_q=\frac1q\sum_{j=0}^{q-1}jT^j.
\]
The identity \((T-1)R_q=1-\Pi\) gives
\(\partial^{-1}=t_0R_qT_{t_0/2}\) on \(\ker\Pi\), of norm at most
\(Cq|t_0|\).  Projected multiplication
\((1-\Pi)a_*:\ker\Pi\to\ker\Pi\) differs from multiplication by
\(1/4\) by at most \(2\|a_*-1/4\|\), so is invertible by a Neumann
series.  Composing these three inverses yields
\[
 \|\{\partial(a_*\partial)\}^{-1}\|_{\ker\Pi}
 \le Cq^2t_0^2.
\]
The projection commutes with all translations; no trivial action of the
half-shift on its range is needed.  We may therefore correct
\(\dot V_*\) by \(V_*'z\), where \(z=O_{\mathrm p}(\eps)\), to
make its total residual jet zero.

Finally set \(w=(V_*-V)/V'\) and differentiate this quotient using the
corrected jet of \(V_*\).  Both \(w\) and \(\dot w\), with the fixed
number of derivatives needed below, are \(O_{\mathrm p}(\eps)\) after
the losses up to \(Y_{10}\).  Taylor's formula and
\eqref{joint:linearization} give
\[
 0=E(V_*)=E+\partial(a\partial w)+(Ew)'
              +O_{\mathrm p}(\eps^2).
\]
The same formula differentiated in the assigned jets has a quadratic
remainder.  To see this directly, write the remainder as the integral
of the second variation along the segment from \(V\) to \(V_*\).
A first parameter derivative either replaces one of its two increments
by its assigned jet, or differentiates a coefficient.  The first case
still contains two \(O_{\mathrm p}(\eps)\) factors.  The second costs
only fixed derivatives in \eqref{joint:raw-jet}, with denominators at
worst \(t_0^{-4}\), and also retains both factors.  Cauchy estimates
on the remaining fixed buffers absorb the required derivatives.
Moreover \((Ew)'\) and its jet are quadratic since
\(E,E_t,w,\dot w=O_{\mathrm p}(\eps)\).  This proves both
\eqref{joint:value-compatibility} and \eqref{joint:jet-compatibility}.
All costs are fixed powers of \(Q\): the mesh inverse costs \(Q^2\),
its Hessian \(Q^3\), the cyclic inverse at most \(Q^2\), and all other
steps use fixed derivatives and fixed Cauchy losses.  Negative \(t_0\)
follow by evenness, with the sign of the first jet reversed.
\end{proof}

\subsection{Two Fourier divisions}

For a periodic function \(F\), write
\(F(y,t)=\sum_{n\in\mathbb Z}F_n(t)e^{\ii ny}\).  The nonzero-mode
divisor of \(\partial\) is
\begin{equation}
 d_n(t)=\frac{2\ii\sin(nt/2)}{t},\qquad d_n(0)=\ii n.
 \label{joint:divisor}
\end{equation}
Its nonzero zeros \(t_j=2\pi j/n\), \(j\ne0\), are simple.
We use the following analytic estimate, with exactly the buffers needed
here.

\begin{lemma}[Buffered division, {\cite[Lemma~3.5]{NB}}]
\label{joint:division}
Let \(\nu\in\{1/2,1\}\), and suppose a mean-zero periodic function
\(F\) is holomorphic on \(\mathcal D_{j_0}^\nu\), with norm at most
\(M\).  Let the fixed integer buffers satisfy
\(b_{\mathrm{buf}}\ge j_0+1\) and \(a_{\mathrm{buf}}\ge j_0\).  Suppose, for
\(1\le |n|\le N\) and every divisor zero with
\(0<|t_j|<d-b_{\mathrm{buf}}\delta\), that
\[
 |F_n(t_j)|\le\eta e^{-|n|(s-a_{\mathrm{buf}}\delta)}.
\]
For the retained modes define
\begin{align}
 R_n(t)&=\sum_{0<|t_j|<d-b_{\mathrm{buf}}\delta} F_n(t_j)
       \frac{\sin(n(t-t_j)/2)}{n(t-t_j)/2},
       \label{joint:interpolation}\\
 H_n(t)&=\frac{F_n(t)-R_n(t)}{d_n(t)},
       \label{joint:quotient}
\end{align}
and set \(H_0=0\) and \(H_n=0\) for \(|n|>N\).
The quotients are holomorphic for \(|t|<d-(b_{\mathrm{buf}}+1)\delta\).
For each fixed integer \(J>\max(a_{\mathrm{buf}},j_0,b_{\mathrm{buf}}+1)\),
\begin{align}
 \|H\|_{\mathcal D_J^{\nu-1/2}}&\le CQ^B(M+\eta),\notag\\
 \|\partial H-F\|_{\mathcal D_J^\nu}
   &\le CQ^B\bigl(\eta+Me^{-N\delta}\bigr).
 \label{joint:division-bounds}
\end{align}
Here \(C,B\) depend on the fixed buffers, not on \(N,d,\delta\).
Fixed additional derivatives cost fixed additional powers after fixed
domain losses.  Real symmetry and evenness are preserved.
\end{lemma}

The interpolation in \eqref{joint:interpolation} cancels precisely the
included zeros, since its sinc kernel is one at its own center and
zero at every other center.  The buffers ensure that every possible
pole in the smaller parameter disk has been canceled.  The gain of
one half in the domain exponent comes from the exponential factor
\(e^{-|n\Im t|/2}\) in the inverse sine divisor.  The estimate also
provides, on the real parameter axis,
\begin{equation}
 |R_n(t)|+N^{-1}|R_n'(t)|
   \le CQ^B\eta e^{-|n|(s-a_{\mathrm{buf}}\delta)}.
 \label{joint:interpolation-jet}
\end{equation}
This follows directly from the sinc formula: there are at most \(CN\)
centers and the kernel and its first real derivative are bounded by
\(C\) and \(C|n|\).  Thus the cited estimate includes all derivative
information used below.

\begin{proof}[Proof of Proposition~\ref{joint:newton}]
The first division takes \(F=-E\).  Its mean vanishes by
Lemma~\ref{joint:factorization}.  Every nonzero divisor zero for a mode
\(|n|\le N\) has reduced denominator at most \(N\).  At those centers
with \(|t_j|<d-3\delta\), Equation~\eqref{joint:value-compatibility}
and \(d_n(t_j)=0\) give
\[
 |E_n(t_j)|\le CQ^B\eps^2e^{-|n|(s-10\delta)}.
\]
Apply Lemma~\ref{joint:division} with
\[
 (j_0,\nu,a_{\mathrm{buf}},b_{\mathrm{buf}},J)=(0,1,10,3,11).
\]
Denote the result by \(H\).  Then
\begin{equation}
 H=O_{\mathrm p}(\eps)\quad\hbox{on }\mathcal D_{11}^{1/2},
 \qquad
 \partial H+E=O_{\mathrm p}(\eps^2)
          \quad\hbox{on }\mathcal D_{11}^{1}.
 \label{joint:first-division}
\end{equation}
The extra polynomial factor multiplying \(\eps^2\) in the bound for
\(H\) is absorbed by the smallness requirement.  Define
\begin{equation}
 b=a^{-1}(H+\lambda),\qquad
 \lambda(t)=-\frac{\langle a^{-1}H\rangle}{\langle a^{-1}\rangle}.
 \label{joint:mean-correction}
\end{equation}
The denominator stays away from zero, so \(b\) has mean zero and is
\(O_{\mathrm p}(\eps)\) on \(\mathcal D_{12}^{1/2}\).

Before dividing again, we need the compatibility of \(b\) at its
resonances.  Fix a retained rational parameter \(t_0\) with
\(|t_0|<d-14\delta\), and let \(g=a\partial w\) use the value and
assigned jet from Lemma~\ref{joint:compatibility}.  For every retained
nonzero mode, the exact first-division identity gives, as an identity
of values and first jets at \(t_0\),
\begin{equation}
 d_n(H_n-g_n)=-(E+\partial g)_n-R_n.
 \label{joint:comparison}
\end{equation}
The right-hand side and its first jet are bounded by
\(CQ^B\eps^2e^{-|n|(s-10\delta)}\), using both compatibility equations
and \eqref{joint:interpolation-jet}.

If \(d_n(t_0)\ne0\), reducedness of \(t_0/(2\pi)=l/q\) gives
\[
 |\sin(\pi nl/q)|\ge c/q,\qquad
 |d_n(t_0)|^{-1}\le Cq|t_0|.
\]
If \(d_n(t_0)=0\), then
\( |d_n'(t_0)|=|n|/|t_0|\), and differentiating
\eqref{joint:comparison} bounds the same difference \(H_n-g_n\).
Thus in both cases
\[
 |H_n(t_0)-g_n(t_0)|
 \le CQ^B\eps^2 e^{-|n|(s-10\delta)}.
\]
This is the reason for the first-jet compatibility: the value equation
alone would give no estimate for a resonant coefficient of \(g\).

The function \(g\) has norm \(O_{\mathrm p}(\eps)\) on \(Y_{10}\).
Its omitted Fourier tail on \(Y_{13}\) is bounded by
\(CQ^B\eps e^{-3N\delta}=O_{\mathrm p}(\eps^2)\).  Summing the
retained estimates on that smaller strip therefore gives
\[
 H-(g-g_0)=O_{\mathrm p}(\eps^2)\quad\hbox{on }Y_{13}.
\]
Since \(\langle a^{-1}g\rangle=\langle\partial w\rangle=0\), the
mean correction in \eqref{joint:mean-correction} satisfies
\(\lambda=g_0+O_{\mathrm p}(\eps^2)\).  Hence
\begin{equation}
 b-\partial w=O_{\mathrm p}(\eps^2)\quad\hbox{on }Y_{13}.
 \label{joint:second-compatibility}
\end{equation}
At a zero of \(d_n\), the coefficient \((\partial w)_n\) vanishes.
Equation~\eqref{joint:second-compatibility} supplies the second division
with its required Fourier values.

Apply Lemma~\ref{joint:division} to \(F=b\), with
\[
 (j_0,\nu,a_{\mathrm{buf}},b_{\mathrm{buf}},J)=(12,1/2,13,15,17).
\]
All its centers satisfy the comparison range \(|t_0|<d-14\delta\).
Its output \(v\) satisfies
\begin{equation}
 v=O_{\mathrm p}(\eps)\quad\hbox{on }\mathcal D_{17}^{0},
 \qquad
 \partial v-b=O_{\mathrm p}(\eps^2)
       \quad\hbox{on }\mathcal D_{17}^{1/2}.
 \label{joint:second-division}
\end{equation}
Both divisions, multiplication by \(a^{-1}\), and the mean adjustment
preserve parity and real symmetry.

Since \(ab=H+\lambda(t)\) and \(\partial\lambda=0\),
\eqref{joint:first-division} and \eqref{joint:second-division} yield
\[
 E+\partial(a\partial v)=O_{\mathrm p}(\eps^2)
       \quad\hbox{on }\mathcal D_{18}^{1}.
\]
Here the remaining difference costs only a derivative and a domain
loss, by
\(\partial f=\int_{-1/2}^{1/2}f'(y+ut,t)\,\dd u\).
Take \(\Delta V=V'v\).  Cauchy estimates between indices 17 and 20
bound its first three \(y\)-derivatives by
\(O_{\mathrm p}(\eps)\).  Equation~\eqref{joint:linearization} makes
its linearized residual quadratic, because
\((Ev)'=O_{\mathrm p}(\eps^2)\).  The nonsingular Formula
\eqref{joint:regular-formula} gives
\[
 E(V+\Delta V)-E(V)-E_V(\Delta V)=O_{\mathrm p}(\eps^2)
       \quad\hbox{on }\mathcal D_{20}^{1},
\]
including \(t=0\).

The construction uses fixed numbers of matrix inversions, contractions,
products and derivatives, two applications of Lemma~\ref{joint:division},
and fixed Cauchy losses.  All dimension and resonance costs are bounded
by powers of \(N\), and all shrinking-domain costs by powers of
\(\delta^{-1}\).  The action bounds are fixed before \(d\) is shrunk.
A single exponent \(B_0\) and constant \(C_0\) thus give
\eqref{joint:step-size}--\eqref{joint:step-error} and absorb every
smallness condition.  This proves the Proposition with the asserted
uniformity.
\end{proof}

\subsection{Initialization and convergence}

For any fixed integer \(m\ge1\), there are analytic periodic
coefficients \(v_1,\ldots,v_m\) such that
\begin{equation}
 V^{(m)}(y,t)=y+\sum_{j=1}^m v_j(y)t^{2j},\qquad
 E(V^{(m)})=O(t^{2m+2}).
 \label{joint:initializer}
\end{equation}
Here is the recursion, also given in \cite[Lemma~3.6]{NB}.
At order \(t^{2j}\), the new coefficient enters as \(v_j''/4\), by
\eqref{joint:at-zero}; all other terms depend on earlier coefficients.
Their mean is zero by the mean-zero identity for \(E\).  Solve the
periodic second-derivative equation, choosing \(v_j\) of mean zero.
Each fixed finite recursion uses only finitely many derivatives of
functions holomorphic on the original open strip.  Thus every such
recursion is valid on one fixed smaller strip, with constants allowed
to depend on \(m\).  Integration preserves holomorphy and real symmetry.
In particular, on that fixed strip and for small \(d_0\),
\[
 \|E(V^{(m)})\|\le C_m d_0^{2m+2},\qquad
 \|V^{(m)}-y\|_{C_y^3}\le C_m d_0^2.
\]
No bound uniform in \(m\) is required.

\begin{proof}[Proof of Theorem~\ref{thm:joint}]
Fix a sufficiently small \(y\)-width \(s_0\); all constants in
Proposition~\ref{joint:newton} are now fixed.  Begin with
\eqref{joint:initializer} on radius \(d_0<s_0/4\), choosing \(m\)
and then \(d_0\) below.  At step \(i\), take
\[
 \delta_i=2^{-i}d_0/200
\]
and reduce both current radii by \(20\delta_i\).  Their limits satisfy
\(d_\infty\ge4d_0/5\) and \(s_\infty\ge s_0-d_0/5\).
If the current residual norm is \(\eps_i>0\), choose
\[
 N_i=\left\lceil\delta_i^{-1}\log(1/\eps_i)\right\rceil.
\]
If it vanishes, an exact solution has already been obtained.  Every
fixed power of \(1+\log(1/\eps_i)\) is bounded by a constant times
\(\eps_i^{-1/4}\).  Thus, with one fixed enlarged exponent \(B\),
Proposition~\ref{joint:newton} gives
\begin{align}
 \eps_{i+1}&\le C d_0^{-B}2^{Bi}\eps_i^{7/4},
 \label{joint:error-iteration}\\
 \|\Delta V_i\|_{C_y^3}
   &\le C d_0^{-B}2^{Bi}\eps_i^{3/4}
 \label{joint:correction-iteration}
\end{align}
on the next domains, provided the polynomial smallness requirement holds.

Choose \(m\) so that \(2m+2>4B+4\), and only then shrink \(d_0\).
The finite constant \(C_m\) is thereby absorbed, giving the initial
bound in the induction
\begin{equation}
 \eps_i\le d_0^{4B+4}2^{-(4B+4)i}.
 \label{joint:majorant}
\end{equation}
Indeed substituting this bound in \eqref{joint:error-iteration} and
dividing by the desired next bound gives at most
\[
 C d_0^{2B+3}2^{4B+4-(2B+3)i},
\]
which is at most one for small \(d_0\).  The correction bound becomes
\[
 \|\Delta V_i\|_{C_y^3}
 \le C d_0^{2B+3}2^{-(2B+3)i}.
\]
These bounds both imply the polynomial smallness condition at every
step and keep the candidates in their fixed smallness class: the
initializer has slack and the total correction is arbitrarily small
as \(d_0\to0\).  The induction is therefore justified before each
application of the Newton estimate.  The order of choices is
\(B\), then \(m\), then \(d_0\).

The corrections converge uniformly, with their first three
\(y\)-derivatives, on the limiting domains.  Their limit is
holomorphic, real-symmetric, even in \(t\), and periodic modulo \(y\),
with \(V_y\) close to one.  The residual tends to zero, so
\eqref{joint:stationary} holds wherever the two shifts remain in the
limit domain.  At \(t=0\), Equation~\eqref{joint:at-zero} gives
\(V''=0\); periodicity gives \(V(y,0)=y+c\), with real \(c\).
Replacing \(y\) by \(y-c\) removes the constant and preserves the
equation and all symmetries.  Finally choose a smaller product
neighborhood with enough \(y\)-margin for both shifts.  This gives
precisely the domain and conclusions stated in the Theorem.
\end{proof}
